\chapter{Esecuzione e risultati}
\label{cap:td-esecuzione}

\section{Esito complessivo}

\begin{table}[H]
\centering
\begin{tabular}{@{}L{7.6cm} r@{}}
\toprule
\textbf{Grandezza} & \textbf{Valore} \\
\midrule
Suite eseguite            & 13 \\
Casi di test eseguiti     & 232 \\
Casi superati             & 232 \\
Casi falliti              & 0 \\
Tempo di esecuzione       & inferiore a 10 secondi \\
\bottomrule
\end{tabular}
\caption{Esito dell'esecuzione della suite automatica.}
\label{tab:td-esito}
\end{table}

Il numero di asserzioni eseguite non coincide con quello dei casi specificati
nel \cref{cap:td-casi} per due ragioni. Alcuni casi sono parametrici: un solo
caso specificato, come \id{TC\_AX\_01}, corrisponde a un'asserzione ripetuta
su ciascuna delle rotte riservate. Le nove prove di sistema, viceversa, sono
manuali e non compaiono fra le asserzioni automatiche. La corrispondenza fra
le due numerazioni è riportata nella tabella seguente.

\begin{table}[H]
\centering
\begin{tabular}{@{}L{6.2cm} r r@{}}
\toprule
\textbf{Area} & \textbf{Casi specificati} & \textbf{Asserzioni eseguite} \\
\midrule
Autenticazione                & 14 & 29 \\
Gestione account              & 10 & 13 \\
Gestione tutorial             & 21 & 28 \\
Gestione quiz                 & 30 & 29 \\
Gestione feedback             &  8 & 17 \\
Gestione obiettivi            & 18 & 19 \\
Serializzazione e validazione &  — & 25 \\
Controllo degli accessi       & 12 & 50 \\
Infrastruttura                & 14 & 22 \\
Prove di sistema (manuali)    &  9 & — \\
\midrule
\textbf{Totale}               & \textbf{136} & \textbf{232} \\
\bottomrule
\end{tabular}
\caption{Casi di test specificati e asserzioni effettivamente eseguite.}
\label{tab:td-conteggi}
\end{table}

\begin{nota}
L'area \enquote{Gestione quiz} è l'unica in cui le asserzioni sono meno dei
casi specificati: tre casi relativi al livello HTTP sono contati nell'area
\enquote{Serializzazione e validazione}, dove risiede la suite che li
contiene. La ripartizione delle asserzioni segue l'organizzazione dei file di
test, quella dei casi segue l'organizzazione funzionale del documento.
\end{nota}

\section{Difetti individuati e risolti}
\label{sec:td-difetti}

I difetti elencati sono stati individuati durante la realizzazione del
sistema, in tre modi distinti: dall'analisi statica e dal controllo dei tipi,
dalla scrittura dei casi di test, e dalle prove di sistema. Tutti sono
chiusi. Sono riportati perché rendono conto di che cosa la verifica ha
effettivamente trovato, e perché la loro distribuzione dice qualcosa su quale
tecnica ha reso di più.

\subsection{Difetti di sicurezza}

\begin{longtable}{@{}l L{5.6cm} L{6.4cm}@{}}
\toprule
\textbf{Id} & \textbf{Difetto} & \textbf{Risoluzione e caso di test} \\
\midrule
\endfirsthead
\toprule
\textbf{Id} & \textbf{Difetto} & \textbf{Risoluzione e caso di test} \\
\midrule
\endhead
\bottomrule
\endlastfoot

\id{DF\_01} &
Le password erano confrontate in chiaro e conservate senza hashing. &
Hashing bcrypt in registrazione e confronto in accesso.
\id{TC\_AU\_04}. \\

\id{DF\_02} &
La risposta all'accesso conteneva la password, che veniva conservata nel
browser. &
Introdotto il livello DTO. \id{TC\_AU\_11}. \\

\id{DF\_03} &
Nessuna rotta verificava identità o ruolo: qualunque chiamante poteva
eliminare contenuti e leggere l'elenco degli utenti. &
Middleware di autenticazione e autorizzazione.
\id{TC\_AX\_01}--\id{TC\_AX\_06}. \\

\id{DF\_04} &
L'autore di feedback e svolgimenti era preso dal corpo della richiesta. &
L'identità è ricavata dalla sessione. \id{TC\_AX\_09}, \id{TC\_AX\_10}. \\

\id{DF\_05} &
Il contenuto dei tutorial era memorizzato senza sanificazione e reso come
HTML. &
Sanificazione in ingresso. \id{TC\_CT\_07}, \id{TC\_CT\_08}. \\

\id{DF\_06} &
L'eliminazione di un'immagine accettava un percorso arbitrario, consentendo
di uscire dalla cartella consentita. &
Solo nomi semplici, risolti e verificati.
\id{TC\_IN\_04}--\id{TC\_IN\_11}. \\

\id{DF\_07} &
Le risposte corrette dei quiz erano inviate al client, che calcolava l'esito. &
Correzione lato server e rimozione del campo dalla rappresentazione.
\id{TC\_QZ\_29}. \\

\end{longtable}

\subsection{Difetti funzionali}

\begin{longtable}{@{}l L{5.6cm} L{6.4cm}@{}}
\toprule
\textbf{Id} & \textbf{Difetto} & \textbf{Risoluzione e caso di test} \\
\midrule
\endfirsthead
\toprule
\textbf{Id} & \textbf{Difetto} & \textbf{Risoluzione e caso di test} \\
\midrule
\endhead
\bottomrule
\endlastfoot

\id{DF\_08} &
Le risposte dell'utente erano associate alle domande per posizione: un ordine
di presentazione diverso produceva una correzione errata in silenzio. &
Associazione per identificativo. \id{TC\_QZ\_18}. \\

\id{DF\_09} &
Il contatore dei quiz superati veniva incrementato a ogni superamento, anche
ripetendo un quiz già superato: i badge venivano assegnati per traguardi mai
raggiunti. &
Il valore è ricalcolato dagli svolgimenti registrati.
\id{TC\_QZ\_21}, \id{TC\_QZ\_23}. \\

\id{DF\_10} &
La soglia di superamento era $> 70\%$ nel client e $\ge 70\%$ nel server: un
punteggio esattamente pari alla soglia produceva esiti discordanti. &
Soglia unica, definita in un solo punto. \id{TC\_QZ\_17}. \\

\id{DF\_11} &
Il limite di lunghezza del commento era 500 caratteri nell'applicazione e 280
nel database: un commento fra i due valori era accettato e poi rifiutato dal
DBMS. &
Limite allineato a 500 su entrambi i livelli.
\id{TC\_FB\_06}, \id{TC\_SY\_06}. \\

\id{DF\_12} &
Costruzione di oggetti \id{Risposta} con argomenti in ordine errato e
riferimenti a colonne inesistenti nel livello di persistenza. &
Livello di persistenza riscritto con tipi di riga espliciti.
Individuato dal controllo dei tipi; coperto da \id{TC\_SY\_03}. \\

\id{DF\_13} &
Il livello di persistenza degli obiettivi passava riferimenti a funzione al
posto dei valori, e alcune istruzioni di aggiornamento avevano i parametri in
ordine invertito. &
Corretti i riferimenti e l'ordine dei parametri. Individuato in revisione del
codice; coperto da \id{TC\_SY\_03}. \\

\id{DF\_14} &
L'eliminazione di un quiz caricava tutte le domande del database e le
filtrava in memoria, pur essendo la cascata già garantita dai vincoli di
integrità. &
Rimossa la cascata applicativa. \id{TC\_QZ\_13}, \id{TC\_SY\_07}. \\

\end{longtable}

\subsection{Difetti individuati dalle prove di sistema}

\begin{longtable}{@{}l L{5.6cm} L{6.4cm}@{}}
\toprule
\textbf{Id} & \textbf{Difetto} & \textbf{Risoluzione e caso di test} \\
\midrule
\endfirsthead
\toprule
\textbf{Id} & \textbf{Difetto} & \textbf{Risoluzione e caso di test} \\
\midrule
\endhead
\bottomrule
\endlastfoot

\id{DF\_15} &
Durante la correzione di un quiz l'aggiornamento del contatore dei quiz
superati usava il pool di connessioni invece della connessione della
transazione in corso: la richiesta attendeva i lock che essa stessa deteneva,
e terminava con un errore di attesa scaduta. &
La connessione transazionale è propagata a tutti i DAO coinvolti.
Individuato da \id{TC\_SY\_03}; coperto da \id{TC\_OB\_05} e da
un'asserzione dedicata sulla propagazione. \\

\end{longtable}

\begin{nota}
\textbf{Che cosa dice la distribuzione dei difetti.} Sette difetti su
quindici riguardano la sicurezza e nessuno di essi sarebbe emerso da test
funzionali: erano tutti casi in cui il sistema faceva correttamente ciò che
gli veniva chiesto, ma da chiunque lo chiedesse. Sono stati trovati
formulando i requisiti non funzionali in modo verificabile e scrivendo i casi
di test corrispondenti.

\id{DF\_15} è l'unico difetto individuato esclusivamente dalle prove di
sistema, e riguarda proprio l'area che la verifica automatica non copre: il
comportamento transazionale reale del DBMS. È la conferma sperimentale del
limite dichiarato in \cref{sec:td-limiti}.
\end{nota}

\section{Copertura del codice}
\label{sec:td-copertura}

\subsection{Perimetro di misurazione}

La copertura è misurata sul codice che contiene decisioni: servizi, rotte,
middleware, errori e serializzazione. Le entità sono escluse dal perimetro:
sono contenitori di dati con soli metodi di accesso, esercitati indirettamente
da tutte le suite, e includerle gonfierebbe la metrica senza aggiungere
garanzie.

\subsection{Risultati}

\begin{longtable}{@{}L{7.6cm} r r r r@{}}
\toprule
\textbf{Modulo} & \textbf{Istr.} & \textbf{Diram.} & \textbf{Funz.} & \textbf{Righe} \\
\midrule
\endfirsthead
\toprule
\textbf{Modulo} & \textbf{Istr.} & \textbf{Diram.} & \textbf{Funz.} & \textbf{Righe} \\
\midrule
\endhead
\bottomrule
\endlastfoot

\file{services/AutenticazioneService} & 100,0 & 100,0 & 100,0 & 100,0 \\
\file{services/AccountService}        &  97,2 & 100,0 & 100,0 &  97,2 \\
\file{services/TutorialService}       & 100,0 &  85,7 & 100,0 & 100,0 \\
\file{services/QuizService}           &  97,8 &  88,9 & 100,0 &  97,8 \\
\file{services/FeedbackService}       & 100,0 &  94,4 & 100,0 & 100,0 \\
\file{services/ObiettivoService}      & 100,0 &  91,7 & 100,0 & 100,0 \\
\midrule
\file{routes/auth}                    & 100,0 & 100,0 & 100,0 & 100,0 \\
\file{routes/accounts}                & 100,0 & 100,0 & 100,0 & 100,0 \\
\file{routes/tutorials}               & 100,0 &  87,5 & 100,0 & 100,0 \\
\file{routes/quiz}                    &  92,3 & 100,0 &  83,3 &  92,3 \\
\file{routes/feedback}                & 100,0 & 100,0 & 100,0 & 100,0 \\
\file{routes/obiettivi}               &  91,3 & 100,0 &  80,0 &  91,3 \\
\midrule
\file{middleware/auth}                &  84,9 &  62,5 &  88,9 &  84,0 \\
\file{middleware/errorHandler}        &  93,8 &  66,7 & 100,0 &  92,9 \\
\file{middleware/validate}            & 100,0 &  85,7 & 100,0 & 100,0 \\
\file{middleware/asyncHandler}        & 100,0 & 100,0 & 100,0 & 100,0 \\
\file{middleware/upload}              &  62,2 &  60,0 &  33,3 &  63,6 \\
\midrule
\file{errors/AppError}                & 100,0 & 100,0 & 100,0 & 100,0 \\
\file{dto/index}                      & 100,0 &  50,0 & 100,0 & 100,0 \\
\midrule
\textbf{Totale} & \textbf{94,9} & \textbf{86,9} & \textbf{93,1} & \textbf{95,0} \\

\end{longtable}

In valore assoluto: 615 istruzioni coperte su 648, 146 diramazioni su 168,
122 funzioni su 131, 605 righe su 637. Le soglie configurate — 85\% di
istruzioni, 80\% di diramazioni, 85\% di funzioni e righe — sono superate, e
la loro violazione farebbe fallire la suite.

\subsection{Analisi delle parti scoperte}

\begin{longtable}{@{}L{4.9cm} L{9.5cm}@{}}
\toprule
\textbf{Modulo} & \textbf{Che cosa resta scoperto e perché} \\
\midrule
\endfirsthead
\toprule
\textbf{Modulo} & \textbf{Che cosa resta scoperto e perché} \\
\midrule
\endhead
\bottomrule
\endlastfoot

\file{middleware/upload} &
È il modulo con la copertura più bassa (62,2\%). La parte scoperta è la
configurazione di multer e il ridimensionamento delle immagini: codice di
inizializzazione e chiamate a librerie esterne, esercitati dalle prove di
sistema \id{TC\_SY\_05}. La logica applicativa del modulo — generazione del
nome, validazione del percorso — è invece interamente coperta, ed è la parte
che ha rilevanza per la sicurezza. \\

\file{middleware/auth} &
Restano scoperti alcuni rami difensivi nella decodifica del token: casi in
cui il payload è formalmente valido ma privo dei campi attesi. Sono
raggiungibili solo con un token firmato dal server e malformato, condizione
che non può verificarsi se non manomettendo la chiave. \\

\file{middleware/errorHandler} &
Il ramo scoperto è la registrazione dell'errore nei log, disattivata in
ambiente di test per non inquinare l'output. \\

\file{dto/index} &
La diramazione scoperta è la conversione di una data assente. Riguarda
soltanto feedback costruiti in memoria e non ancora persistiti. \\

\file{routes/quiz}, \file{routes/obiettivi} &
Restano scoperte le rotte di aggiornamento, il cui comportamento è identico a
quello delle rotte di creazione già coperte e delega interamente al servizio,
a sua volta coperto. \\

\file{services/TutorialService} &
Le diramazioni scoperte sono i rami di messaggio d'errore alternativi di uno
stesso controllo di validazione, distinguibili solo dal testo prodotto. \\

\end{longtable}
